\documentclass{article}

% packages
\usepackage{amsmath, amsthm, thmtools, amsfonts, amssymb, luacode, catchfile, tikzducks, hyperref, ifthen}
\ifcsname c@kobocompile\endcsname
	\usepackage[a5paper, total={1072pt, 1448pt}, margin=10pt, includeheadfoot]{geometry} % set page margins
\else
	\usepackage[a4paper, margin=50pt, includeheadfoot]{geometry}
\fi
\usepackage[shortlabels]{enumitem}
\usepackage[skip=3pt, indent=0pt]{parskip}

% language
\usepackage[bidi=basic, layout=tabular, provide=*]{babel}
\ifcsname c@english\endcsname
	\babelprovide[main, import]{english}
\else
	\babelprovide[main, import]{hebrew}
	\babelprovide{rl}
\fi
%\babelfont{rm}{Libertinus Serif}
\babelfont{rm}[Renderer=Harfbuzz]{Libertinus Serif}
\babelfont{sf}{Libertinus Sans}
\babelfont{tt}{Libertinus Mono}

% style
\AddToHook{cmd/section/before}{\clearpage}	% Add line break before section
\linespread{1.3}
\setcounter{secnumdepth}{0}		% Remove default number tags from sections, this won't do well with theorems
\AtBeginDocument{\setlength{\belowdisplayskip}{3pt}}
\AtBeginDocument{\setlength{\abovedisplayskip}{3pt}}
\graphicspath{ {../images/} }

% operators
\DeclareMathOperator\cis{cis}
\DeclareMathOperator\Sp{Sp}
\DeclareMathOperator\tr{tr}
\DeclareMathOperator\im{Im}
\DeclareMathOperator\re{Re}
\DeclareMathOperator\diag{diag}
\DeclareMathOperator*\lowlim{\underline{lim}}
\DeclareMathOperator*\uplim{\overline{lim}}
\DeclareMathOperator\rng{rng}
\DeclareMathOperator\Sym{Sym}
\DeclareMathOperator\Arg{Arg}
\DeclareMathOperator\Log{Log}
\DeclareMathOperator\dom{dom}
\DeclareMathOperator\supp{Supp}
\DeclareMathOperator\var{Var}
\DeclareMathOperator\cov{Cov}

% commands
%\renewcommand\qedsymbol{\textbf{מש''ל}}
%\renewcommand\qedsymbol{\fbox{\emoji{lizard}}}
\newcommand{\Aa}[0]{\mathcal{A}}
\newcommand{\Bb}[0]{\mathcal{B}}
\newcommand{\CC}[0]{\mathbb{C}}
\newcommand{\Cc}[0]{\mathcal{C}}
\newcommand{\EE}[0]{\mathbb{E}}
\newcommand{\FF}[0]{\mathbb{F}}
\newcommand{\Ff}[0]{\mathcal{F}}
\newcommand{\Ii}[0]{\mathcal{I}}
\newcommand{\Gg}[0]{\mathcal{G}}
\newcommand{\Ll}[0]{\mathcal{L}}
\newcommand{\Mm}[0]{\mathcal{M}}
\newcommand{\NN}[0]{\mathbb{N}}
\newcommand{\Nn}[0]{\mathcal{N}}
\newcommand{\PP}[0]{\mathbb{P}}
\newcommand{\Pp}[0]{\mathcal{P}}
\newcommand{\QQ}[0]{\mathbb{Q}}
\newcommand{\RR}[0]{\mathbb{R}}
\newcommand{\Rr}[0]{\mathcal{R}}
\newcommand{\Ss}[0]{\mathcal{S}}
\newcommand{\TT}[0]{\mathbb{T}}
\newcommand{\Uu}[0]{\mathcal{U}}
\newcommand{\Vv}[0]{\mathcal{V}}
\newcommand{\Ww}[0]{\mathcal{W}}
\newcommand{\ZZ}[0]{\mathbb{Z}}
\newcommand{\acts}[0]{\circlearrowright}
\newcommand{\explain}[2] {
	\begin{flalign*}
		 && \text{#2} && \text{#1}
	\end{flalign*}
}
\newcommand{\maketitleprint}[0]{ \begin{center}
	%\begin{tikzpicture}[scale=3]
	%	\duck[graduate=gray!20!black, tassel=red!70!black]
	%\end{tikzpicture}	
	\includegraphics[width=6cm]{cover}
\end{center}
}

% theorem commands
\newtheoremstyle{c_remark}
	{}	% Space above
	{}	% Space below
	{}% Body font
	{}	% Indent amount
	{\bfseries}	% Theorem head font
	{}	% Punctuation after theorem head
	{.5em}	% Space after theorem head
	{\thmname{#1}\thmnumber{ #2}\thmnote{ \normalfont{\text{(#3)}}}}	% head content
\newtheoremstyle{c_definition}
	{3pt}	% Space above
	{3pt}	% Space below
	{}% Body font
	{}	% Indent amount
	{\bfseries}	% Theorem head font
	{}	% Punctuation after theorem head
	{.5em}	% Space after theorem head
	{\thmname{#1}\thmnumber{ #2}\thmnote{ \normalfont{\text{(#3)}}}}	% head content
\newtheoremstyle{c_plain}
	{3pt}	% Space above
	{3pt}	% Space below
	{\itshape}% Body font
	{}	% Indent amount
	{\bfseries}	% Theorem head font
	{}	% Punctuation after theorem head
	{.5em}	% Space after theorem head
	{\thmname{#1}\thmnumber{ #2}\thmnote{ \text{(#3)}}}	% head content

\ifcsname c@english\endcsname
	\theoremstyle{plain}
	\newtheorem{theorem}{Theorem}[section]
	\newtheorem{lemma}[theorem]{Lemma}
	\newtheorem{proposition}[theorem]{Proposition}
	\newtheorem*{proposition*}{Proposition}
	%\newtheorem{corollary}[theorem]{אין חלופה עברית}

	\theoremstyle{definition}
	\newtheorem{definition}[theorem]{Definition}
	\newtheorem*{definition*}{Definition}
	\newtheorem{example}{Example}[section]
	\newtheorem{exercise}{Exercise}[section]

	\theoremstyle{remark}
	\newtheorem*{remark}{Remark}
	\newtheorem*{solution}{Solution}
	\newtheorem{conclusion}[theorem]{Conclusion}
	\newtheorem{notation}[theorem]{Notation}
\else
	\theoremstyle{c_plain}
	\newtheorem{theorem}{משפט}[section]
	\newtheorem{lemma}[theorem]{למה}
	\newtheorem{proposition}[theorem]{טענה}
	\newtheorem*{proposition*}{טענה}
	%\newtheorem{corollary}[theorem]{אין חלופה עברית}

	\theoremstyle{c_definition}
	\newtheorem{definition}[theorem]{הגדרה}
	\newtheorem*{definition*}{הגדרה}
	\newtheorem{example}{דוגמה}[section]
	\newtheorem{exercise}{תרגיל}[section]

	\theoremstyle{c_remark}
	\newtheorem*{remark}{הערה}
	\newtheorem*{solution}{פתרון}
	\newtheorem{conclusion}[theorem]{מסקנה}
	\newtheorem{notation}[theorem]{סימון}
\fi

% Questions related commands
\newcounter{question}
\setcounter{question}{1}
\newcounter{sub_question}
\setcounter{sub_question}{1}

\ifcsname c@english\endcsname
	\newcommand{\question}[1][0]{
		\ifthenelse{#1 = 0}{}{\setcounter{question}{#1}}
		\section{Question \arabic{question}}
		\addtocounter{question}{1}
		\setcounter{sub_question}{1}
	}

	\newcommand{\subquestion}[1][0]{
		\ifthenelse{#1 = 0}{}{\setcounter{sub_question}{#1}}
		\subsection{Part \alph{sub_question}}
		\addtocounter{sub_question}{1}
	}
\else
	\newcommand{\question}[1][0]{
		\ifthenelse{#1 = 0}{}{\setcounter{question}{#1}}
		\section{שאלה \arabic{question}}
		\addtocounter{question}{1}
		\setcounter{sub_question}{1}
	}

	\newcommand{\subquestion}[1][0]{
		\ifthenelse{#1 = 0}{}{\setcounter{sub_question}{#1}}
		\subsection{סעיף \localecounter{letters.gershayim}{sub_question}}
		\addtocounter{sub_question}{1}
	}
\fi

% import lua and start of document
\directlua{common = require ('../common')}

\GetEnv{AUTHOR}

% headers
\author{\AUTHOR}
\date\today

\title{פתרון מטלה 03 --- אנליזה על יריעות, 80426}

\begin{document}
\maketitle
\maketitleprint{}

\question{}
יהי $\vec{F} : U \to \RR^2$ עבור $U \subseteq \RR^2$ שדה וקטורי.
נניח שלכל מלבן $R \subseteq U$ האינטגרל בין שני הקצוות שווה לאינטגרל בין שני הקצוות האחרים.
נראה ש־$\vec{F}$ משמר.
\begin{proof}
	האסטרטגיה שלנו היא להשתמש בחלוקה למלבנים לכל מסילה כללית $\gamma : [a, b] \to U$ במטרה להראות שהאינטגרל על המסילה תלוי רק בנקודות קצה.

	תהיינה $\gamma_1, \gamma_2 : [a, b] \to U$ מסילות כלשהן כך שהן חולקות נקודות קצה.
	עבור $N \in \NN$ נגדיר את הסדרה ${\{ x_n \}}_{n = 0}^N \subseteq U$ כך ש־$x_n = \gamma_1(a + (b - a)\frac{n}{N})$, כלומר זוהי חלוקה של המסילה.
	נבנה סדרת מלבנים כך שבין $x_i, x_{i + 1}$ נמצא מלבן לכל $i$, כך שנקודות אלה הן אלכסון המלבן, מטעמי נוחות נבחר את האלכסון הנותר כך שלכל המלבנים תהיה אוריינטציה זהה.
	נניח גם ש־$ABCD$ מלבן שאלכסונו הוא $\gamma_1(a), \gamma_1(b)$ עם אוריינטציה זהה לזו של המלבנים בחלוקה.
	לבסוף נאריך את המלבנים בסדרה כך שצלע שלהם תשב על המלבן $ABCD$.
	מהנתון נובע שאם $B_i$ הקודקודים של המלבנים בסדרה כפי שבחרנו, אז,
	\[
		\int_{\gamma_{ADC}} \vec{F}\ dl
		= \sum_{i = 0}^N \int_{\gamma_{x_i B_i x_{i + 1}}} \vec{F}\ dl
	\]
	ונקבל,
	\[
		\int_{\gamma_{ADC}} \vec{F}\ dl
		\lim_{N \to \infty} \int_{\gamma_{ADC}} \vec{F}\ dl
		= \lim_{N \to \infty} \sum_{i = 0}^N \int_{\gamma_{x_i B_i x_{i + 1}}} \vec{F}\ dl
		= \int_{\gamma_1} \vec{F}\ dl
	\]
	אבל גם,
	\[
		\int_{\gamma_{ADC}} \vec{F}\ dl
		= \int_{\gamma_2} \vec{F}\ dl
	\]
	משיקולים ודמים, ולכן נובע ש־$\vec{F}$ משמר.
\end{proof}

\question{}
נניח ש־$\vec{F}(x, y, z) = (\frac{x}{x^2 + y^2}, \frac{y}{x^2 + y^2}, z)$, בתחום $\RR^3 \setminus \{ (0, 0, z) \}$.

\subquestion{}
נראה ש־$\vec{F}$ משמר.
\begin{proof}
	אילו קיימת $\varphi : \RR^3 \to \RR$ המעידה על שימור $\vec{F}$, אז מתקיים,
	\[
		\frac{\partial}{\partial x} \varphi = \frac{x}{x^2 + y^2}
		\implies \varphi = \frac{1}{2} \ln(x^2 + y^2) + \varphi_{yz}
	\]
	באותו אופן נסיק כי גם,
	\[
		\frac{\partial}{\partial y} \varphi = \frac{y}{x^2 + y^2}
		\implies \varphi = \frac{1}{2} \ln(x^2 + y^2) + \varphi_{xz}
	\]
	ולבסוף שגם,
	\[
		\frac{\partial}{\partial z} \varphi = z
		\implies \varphi = \frac{1}{2} z^2 + \varphi_{xy}
	\]
	נבחין אם כך שאם נגדיר $\varphi(x, y, z) = \frac{1}{2}(\ln(x^2 + y^2) + z^2)$
	אז נקבל בתחום ש־$\varphi$ היא פונקציה רציפה וגזירה, וכן ש־$\nabla \varphi = \vec{F}$, ולכן מהתנאי השקול לשימור נובע ש־$\vec{F}$ אכן משמר.
\end{proof}

\subquestion{}
נחשב את $\int_{\gamma} \vec{F}\ dl$ עבור $\gamma(t) = (\cos t, \sin t, t^2)$ כאשר $t \in [0, 2 \pi]$.
\begin{solution}
	נבחין כי התחום של המסילה הוא לא פשוט קשר, ולכן עלינו לחשב את האינטגרל ישירות מהגדרה,
	\[
		\int_{\gamma} \vec{F}\ dl
		= \int_{0}^{2 \pi} \vec{F}(\gamma(t)) \cdot \gamma'(t)\ dt
		= \int_{0}^{2 \pi} (\frac{\cos t}{1}, \frac{\sin t}{1}, t^2) \cdot (-\sin t, \cos t, 2t)\ dt
		= \int_{0}^{2 \pi} 0 + 2t^2\ dt
		= {\left. \frac{2}{3} t^3 \right|}_{t = 0}^{t = 2 \pi}
		= \frac{2^4 \pi}{3}
	\]
\end{solution}

\question{}
נניח ש־$U \subseteq \RR^n$ תחום פתוח.

\subquestion[2]
נראה שמחלקות הקשירות המסילתית של $U$ הן קבוצות פתוחות, וכן שכל רכיב קשירות מסילתית הוא קשיר מסילתית.
\begin{proof}
	נניח ש־${\{ A_i \}}_{i \in I} \subseteq U$ הן מחלקות הקשירות המסילתית, אז בהכרח $U = \biguplus_{i \in I} A_i$, כלומר זוהי חלוקה של $U$.
	אילו נניח שקיימת $A_i$ לא פתוחה נקבל סתירה ישירה לפתיחות של $U$, לכן נסיק ש־$A_i$ פתוחה לכל $i \in I$.

	נניח עתה ש־$A \subseteq U$ הוא רכיב קשירות מסילתית, ונראה שהוא קשיר מסילתית.
	לכל $x, y \in A$ מתקיים $x \sim y$, כלומר יש מסילה ביניהם, ולכן מצאנו שהקבוצה $A$ קשירה מסילתית.
\end{proof}

\subquestion{}
נראה שמחלקות הקשירות של $U$ הן קבוצות פתוחות, וכן שכל רכיב קשירות הוא קשיר.
\begin{proof}
	החלק הראשון זהה בהוכחתו לזה שבסעיף הקודם.

	נניח ש־$A \subseteq U$ רכיב קשירות.
	בהוכחה של יחס שקילות בתרגול ראינו כי אם $A, B$ קבוצות קשירות כך שהן לא זרות, אז איחודן קשיר.
	נקבע אם כך $x \in A$, ונגדיר גם $B_y \subseteq A$ הקבוצה הקשירה המעידה על $x \simeq y$ לכל $y \in A$.
	נבחין כי בהכרח $U = \bigcup_{y \in A} B_y$, זאת שכן כל נקודה של $U$ מופיעה באיחוד, וכן אילו יש נקודה באיחוד שלא ב־$A$ אז קיבלנו סתירה להיותה רכיב קשירות.
	נקבל מהטענה הקודמת ש־$U$ היא איחוד לא זר של קבוצות קשירות ולכן קשירה.
\end{proof}

\subquestion{}
נראה שקבוצת מחלקות הקשירות (המסילתית) של $U$ היא בת־מניה.
\begin{proof}
	נניח ש־${\{ A_i \}}_{i \in I}$ רכיבי הקשירות (המסילתית), אז לכל $i \in I$ גם $A_i$ היא פתוחה וקשירה (מסילתית).
	יהי $\epsilon > 0$ כך ש־$\operatorname{diam}(A_i) < \epsilon$ לכל $i \in I$.
	נניח בשלילה ש־$|I| > \aleph_0$, ולכן נובע ש־$B_{\epsilon}(x_i)$ עבור כדורים המוכלים במחלקות הקשירות (המסילתית) היא בעצמה קבוצה של כדורים זרים, אבל ידוע כי $\epsilon$־כיסוי של $\RR^n$ הוא בן־מניה בסתירה.
\end{proof}

\question{}
עבור $a, b > 0$ תהי $E_{a, b} = \{ {(x / a)}^2 + {(y / b)}^2 = 1 \} \subseteq \RR^2$ האליפסה עם הפרמטרים $a, b$.
נשתמש במשפט גרין לתחומים פשוטים כדי להראות שמתקיים $\operatorname{Area}(E_{a, b}) = \pi a b$.
\begin{solution}
	נגדיר את שתי הפונקציות $f, g : [-a, a] \to \RR$ על־ידי,
	\[
		f(x) = b \sqrt{1 - {\left(\frac{x}{a}\right)}^2},
		\quad
		g(x) = -f(x)
	\]
	נבחין כי $\rng \gamma_f \cup \rng \gamma_g = \partial E_{a, b}$, כלומר תנאי משפט גרין חלים עבור פונקציות אלה, כאשר $f(\pm a) = g(\pm a)$ ולכן אין דרישת הגדרת ישרים חוסמים מוספים.
	נעבור לחישוב השטח, נגדיר $\vec{F}(x, y) = (0, x)$ ונסיק ממשפט גרין שמתקיים.
	נבחין כי $\gamma_f(t) = (t, b \sqrt{1 - \frac{t^2}{a^2}})$,
	ולכן גם $\vec{F}(\gamma_f(t)) = (0, t)$.
	מהצד השני מתקיים $\gamma_f'(t) = (1, \frac{-b t}{a^2 \sqrt{1 - \frac{t^2}{a^2}}})$
	\begin{align*}
		\operatorname{Area}(E_{a, b})
		& = \int_{\gamma_f} \vec{F}\ dl - \int_{\gamma_g} \vec{F}\ dl \\
		& = 2 \frac{b}{a^2} \cdot \int_{-a}^{a} \frac{-t^2}{\sqrt{1 - {\left(\frac{t}{a}\right)}^2}}\ dt \\
		& = b {\left. t \sqrt{1 - \frac{t^2}{a^2}} - \arcsin\left(\frac{t}{a}\right) \right|}_{t = -a}^{t = a} \\
		& = ab \pi \\
	\end{align*}
\end{solution}

\end{document}
